\section{梯度下降与梯度计算}

\subsection{随机梯度下降（SGD）回顾}

神经网络的训练依赖于\textbf{随机梯度下降}（Stochastic Gradient Descent）：

\begin{figure}[H]
\centering
\includegraphics[width=0.95\textwidth]{slides-images/slide-008.jpg}
\caption{SGD 更新规则：$\theta^{\text{new}} = \theta^{\text{old}} - \alpha \nabla_\theta J(\theta)$，其中 $\alpha$ 是学习率\protect\footnotemark}
\end{figure}
\footnotetext{Slides 第9页。}

$$\theta^{\text{new}} = \theta^{\text{old}} - \alpha \nabla_\theta J(\theta)$$

其中：
\begin{itemize}
    \item $\theta$：模型的所有参数（包括权重矩阵 $W$、偏置 $b$、词向量等）
    \item $\alpha$：学习率（step size）
    \item $\nabla_\theta J(\theta)$：损失函数关于参数的梯度
\end{itemize}

核心问题：\textbf{如何计算 $\nabla_\theta J(\theta)$？} 有两种方法：
\begin{enumerate}
    \item \textbf{手动推导}（By hand）——矩阵微积分
    \item \textbf{算法自动计算}（Algorithmically）——反向传播算法
\end{enumerate}

本讲两种方法都会讲解。

\subsection{本章小结}

SGD 是训练神经网络的核心算法，其关键步骤是计算损失函数关于所有参数的梯度。在深度学习中，参数 $\theta$ 不仅包括传统的权重和偏置，还包括词嵌入等数据表示参数。

%% ========== Part 5: 矩阵微积分 ==========

